% AUTO-GENERATED by the update_record tool. DO NOT EDIT.
% verify_paper regenerates this file from record.json and the run outputs
% and fails on any difference, so manual edits always surface.
\noindent\textbf{Hypothesis 1.} 2026 年 9 月時点で Vercel AI Gateway から利用できる各社の最新モデル 6 種（OpenAI の gpt-6-sol と gpt-6-luna、Anthropic の claude-opus-5.5 と claude-sonnet-5、Google の gemini-3.1-pro-preview と gemini-3.8-flash）は、SciGym 公開コードの Controller をそのまま用いて論文と同じ条件（SciGym-small 137 件、最大 20 反復、temperature 1.0）で解かせたとき、いずれも平均 Simulation Trajectory Error（STE）が論文 Table 1 の最良行（Gemini-2.5-Pro の 0.3212）以下である。~\cite[s1.p1, s1.p5, s1.p10]{duan-2025-measuring}, \cite[s3.p2]{scigymtable1reproduction-2026-scigym}, \cite[s4.p2]{scigymrikyuopenmodels-2026-scigym}
\begin{enumerate}
\item[\textbf{Claim 1}] gpt-6-sol（openai/gpt-6-sol）で SciGym-small 137 件を最大 20 反復で解かせたときの平均 STE は、論文 Table 1 の最良行（Gemini-2.5-Pro）の 0.3212 以下である。~\cite[s1.p2, s1.p3, s1.p4, s1.p5, s1.p10]{duan-2025-measuring}, \cite[s4.p2]{scigymrikyuopenmodels-2026-scigym}, \cite[s2.p1]{h4duan-2025-scigym}
  \emph{Rationale:} OpenAI の現行世代の推論モデルが論文の最良行に届けば、h1 の gpt-6-sol の部分を担う。
  \emph{Criterion:} \texttt{\detokenize{gpt-6-sol}}.\texttt{\detokenize{trajectory_smape}} $-$ 0.3212 $\leq 0$.
  \emph{Prediction:} $[-0.2, 0]$ (open-weight の kimi-k2.6 / kimi-k3 / deepseek-v4.1-flash が同条件で 0.21〜0.24 に達した（s4.p2〜s4.p4）。商用の最新モデルはそれと同等かそれ以上を予測する).
  \emph{Observed:} pending.
  \emph{Verdict:} pending.
\item[\textbf{Claim 2}] gpt-6-luna（openai/gpt-6-luna）で SciGym-small 137 件を最大 20 反復で解かせたときの平均 STE は、論文 Table 1 の最良行（Gemini-2.5-Pro）の 0.3212 以下である。~\cite[s1.p2, s1.p3, s1.p4, s1.p5, s1.p10]{duan-2025-measuring}, \cite[s4.p5]{scigymrikyuopenmodels-2026-scigym}, \cite[s2.p1]{h4duan-2025-scigym}
  \emph{Rationale:} OpenAI の最小クラスのモデルが論文の最良行に届けば、h1 の gpt-6-luna の部分を担う。
  \emph{Criterion:} \texttt{\detokenize{gpt-6-luna}}.\texttt{\detokenize{trajectory_smape}} $-$ 0.3212 $\leq 0$.
  \emph{Prediction:} $[-0.15, 0.05]$ (open-weight の 27B〜35B 級（qwen3.8-27b 0.25、qwen3.6-35b 0.46）が同条件で最良行の前後に分かれた（s4.p5、s4.p6）。最小クラスの商用モデルはその中間を予測する).
  \emph{Observed:} pending.
  \emph{Verdict:} pending.
\item[\textbf{Claim 3}] claude-opus-5.5（anthropic/claude-opus-5.5）で SciGym-small 137 件を最大 20 反復で解かせたときの平均 STE は、論文 Table 1 の最良行（Gemini-2.5-Pro）の 0.3212 以下である。~\cite[s1.p2, s1.p3, s1.p4, s1.p5, s1.p10]{duan-2025-measuring}, \cite[s4.p3]{scigymrikyuopenmodels-2026-scigym}, \cite[s2.p1]{h4duan-2025-scigym}
  \emph{Rationale:} Anthropic の現行の上位モデルが論文の最良行に届けば、h1 の claude-opus-5.5 の部分を担う。
  \emph{Criterion:} \texttt{\detokenize{claude-opus-5.5}}.\texttt{\detokenize{trajectory_smape}} $-$ 0.3212 $\leq 0$.
  \emph{Prediction:} $[-0.2, 0]$ (c1 と同じ（s4.p2〜s4.p4）).
  \emph{Observed:} pending.
  \emph{Verdict:} pending.
\item[\textbf{Claim 4}] claude-sonnet-5（anthropic/claude-sonnet-5）で SciGym-small 137 件を最大 20 反復で解かせたときの平均 STE は、論文 Table 1 の最良行（Gemini-2.5-Pro）の 0.3212 以下である。~\cite[s1.p2, s1.p3, s1.p4, s1.p5, s1.p10]{duan-2025-measuring}, \cite[s4.p4]{scigymrikyuopenmodels-2026-scigym}, \cite[s2.p1]{h4duan-2025-scigym}
  \emph{Rationale:} Anthropic の中位モデルが論文の最良行に届けば、h1 の claude-sonnet-5 の部分を担う。
  \emph{Criterion:} \texttt{\detokenize{claude-sonnet-5}}.\texttt{\detokenize{trajectory_smape}} $-$ 0.3212 $\leq 0$.
  \emph{Prediction:} $[-0.2, 0]$ (中位の商用モデルは open-weight の上位（0.21〜0.25）に近いと予測する（s4.p2〜s4.p5）).
  \emph{Observed:} pending.
  \emph{Verdict:} pending.
\item[\textbf{Claim 5}] gemini-3.1-pro-preview（google/gemini-3.1-pro-preview）で SciGym-small 137 件を最大 20 反復で解かせたときの平均 STE は、論文 Table 1 の最良行（Gemini-2.5-Pro）の 0.3212 以下である。~\cite[s1.p2, s1.p3, s1.p4, s1.p5, s1.p7, s1.p10]{duan-2025-measuring}, \cite[s3.p2]{scigymtable1reproduction-2026-scigym}, \cite[s2.p1]{h4duan-2025-scigym}
  \emph{Rationale:} 論文で最良だった Gemini Pro 系の現行版が自らの旧版の値に届けば、h1 の gemini-3.1-pro-preview の部分を担う。
  \emph{Criterion:} \texttt{\detokenize{gemini-3.1-pro}}.\texttt{\detokenize{trajectory_smape}} $-$ 0.3212 $\leq 0$.
  \emph{Prediction:} $[-0.2, 0]$ (c1 と同じ。なお現行 API の gemini-2.5-pro の再実行は 0.4148 で論文値を再現しなかった（s3.p2）ため、旧版との差は環境差を含む).
  \emph{Observed:} pending.
  \emph{Verdict:} pending.
\item[\textbf{Claim 6}] gemini-3.8-flash（google/gemini-3.8-flash）で SciGym-small 137 件を最大 20 反復で解かせたときの平均 STE は、論文 Table 1 の最良行（Gemini-2.5-Pro）の 0.3212 以下である。~\cite[s1.p2, s1.p3, s1.p4, s1.p5, s1.p10]{duan-2025-measuring}, \cite[s3.p3]{scigymtable1reproduction-2026-scigym}, \cite[s2.p1]{h4duan-2025-scigym}
  \emph{Rationale:} Gemini Flash 系の現行版が論文の最良行に届けば、h1 の gemini-3.8-flash の部分を担う。
  \emph{Criterion:} \texttt{\detokenize{gemini-3.8-flash}}.\texttt{\detokenize{trajectory_smape}} $-$ 0.3212 $\leq 0$.
  \emph{Prediction:} $[-0.15, 0.05]$ (c2 と同じ。現行 API の gemini-2.5-flash の再実行は 0.4465 だった（s3.p3）が、世代が進んだ分を見込む).
  \emph{Observed:} pending.
  \emph{Verdict:} pending.
\end{enumerate}
\noindent\emph{Assumed, so that the claims together imply Hypothesis 1:}
\begin{itemize}
\item temperature 1.0 での 1 回実行の平均 STE（137 件平均）の実行間変動が判定に対して十分小さいこと。先行する Table 1 再現では件ごとの STE の標準偏差はおよそ 0.27、137 件平均の標準誤差はおよそ 0.02 であった（c1〜c6）
\item Vercel AI Gateway 経由の各モデルが、提供元の API から直接呼んだ場合と同じ能力で応答すること。reasoning の強さは各モデルの既定値のまま変えない（c1〜c6）
\item 論文の 6 モデルは 2025 年春の版であり、本研究の 6 モデルは 2026 年 9 月時点の各社の最新版であること。世代の差を測るのが目的なので、同名モデルの再現は求めない（c1〜c6）
\item aarch64 向けの libroadrunner 2.7.0 と antimony 2.14.0 が、論文環境の libroadrunner 2.8.0 と同じ軌道と評価値を与えること（c1〜c6）
\item 平均 STE が「論文の最良行に届くか」を代表すること。RMS（modifier あり / なし）と NTS は同じ実行から表として報告するが、判定基準には用いない（c1〜c6）
\item ゲートウェイの予算超過（HTTP 402）で run が途中で止まらないこと。止まった場合はその run を最初からやり直す（c1〜c6）
\end{itemize}
